\chapter{Testy i~wyniki}
\label{chap:test_and_result}
Eksperymenty są konieczne by wyznaczyć
optymalne parametry dla algorytmu genetycznego,
dodatkowo konieczne jest zbadanie czy rozwiązanie daje
lepsze efekty niż stochastyczne przeszukiwanie przestrzeni.
Poniższe eksperymenty będą używały wartości
błędu \gls{MSE} jako elementu porównawczego.
Można go użyć przy problematyce
klasyfikacyjnej oraz regresyjnej, dzięki temu,
że wyjście w~danych z~wyjściem klasyfikacyjnym musi
być w~formacie $1$~z~$n$ (ang. \textit{one-hot}).
Taki format powoduje, że jest ``równa odległość''
dla każdej niepoprawnej odpowiedzi.

\section{Procedura wyznaczenia optymalnych parametrów}
Parametry, które zostają poddane wyznaczeniu,
to liczba osobników w~populacji i~generacji,
współczynnik w~funkcji korygującej
oraz współczynnik krzyżowania i~mutacji.

Wartość parametru wybierano na podstawie
najniższej średniej wartości \gls{MSE}.

\tocless\subsection{Wyniki}

\section{Procedura badawcza i~ocena skuteczności}
Aby poprawnie zweryfikować skuteczność
rozwiązania przedstawionego w~niniejszej
pracy zostanie użyty eksperyment porównawczy.
Jako punkt odniesienia (ang. \textit{baseline})
zostanie użyte \gls{RS}.
Pozwoli ona wykazać, że zaproponowana metoda optymalizacji
pozwala na znalezienie rozwiązania o~jakości wyższej
niż stochastyczne przeszukiwanie przestrzeni rozwiązań.   

\subsection{Budżet eksperymentów}
W~celu zagwarantowania równych warunków
obu algorytmom konieczne było przydzielenie
dla nich takich samych warunków.
Ze względu na ilość obliczeń koniecznych do wykonania przy
najbardziej optymalnym z~wyznaczonych zestawie parametrów,
został wybrany mniej wymagający obliczeniowo
budżet kosztem trochę mniej optymalnego zbioru.
Zostały one określone przez liczbę potencjalnych rozwiązań,
która została ustalona na $500$ osobników.

Biblioteka \gls{GAAML} wykorzystuje określony budżet poprzez
populację liczącą $50$ osobników oraz liczbę $10$ pokoleń,
co daje $50 \times 10 = 500$ architektur.

Natomiast w~przypadku \gls{RS} zostanie wylosowanych
i~ocenionych dokładnie $500$ niezależnych architektur.

Obie metody będą uczone i~oceniane według
opisu z~\cref{sec:fittness} dla przypadku,
gdzie dane walidacyjne są podane.

\subsection{Metryka}
Ostatnim etapem jest weryfikacja,
który z~sposobów przeszukiwania przestrzeni
rozwiązań okazał się skuteczniejszy.
Zostanie do tego wykorzystany zbiór testowy.

Porównanie będzie odbywało się przy pomocy \gls{MSE}.
Funkcja ta daje wartości ze zbioru $\left[0, \infty\right)$.
Ponieważ określa ona ``odległość''
jednych danych od drugich, to znaczy,
że im większa będzie wartość zwrócona,
tym dane te będą bardziej niedopasowane do oczekiwanych,
a~wartość $0$ oznacza idealne dopasowanie.
Oznacza to, że metoda, która będzie miała mniejszą zwróconą
wartość \gls{MSE} dla najlepszej sieci z~danej metody na
zbiorze testowym, jest przydatniejsza i~skuteczniejsza.

\section{Przeprowadzenie eksperymentów}

\section{Wyniki eksperymentów}
